% Executive Summary -- standalone two-column brief.
% Compile: pdflatex executive_summary.tex
\documentclass[10pt,a4paper,twocolumn]{article}

\usepackage[T1]{fontenc}
\usepackage[utf8]{inputenc}
\usepackage[a4paper,top=1.35cm,bottom=1.5cm,left=1.45cm,right=1.45cm,columnsep=0.65cm]{geometry}
\usepackage{caption}
\usepackage{graphicx}
\usepackage{microtype}
\usepackage{xcolor}
\usepackage{titlesec}

\definecolor{rulegray}{HTML}{555555}

\pagestyle{plain}
\raggedbottom
\setlength{\parindent}{0pt}
\setlength{\parskip}{0.24em}
\titlespacing*{\section}{0pt}{0.65em}{0.25em}
\titleformat{\section}{\large\bfseries}{}{0pt}{}
\captionsetup{font=footnotesize,labelfont=bf,skip=2pt}

\newcommand{\summaryfigure}[4][0.25em]{%
  \par\smallskip
  \noindent\begin{minipage}{\linewidth}
    \centering
    \includegraphics[width=#2\linewidth]{#3}
    \captionof{figure}{#4}
  \end{minipage}
  \par\vspace{#1}
}

\begin{document}
\normalsize

\twocolumn[
\begin{center}
{\large\bfseries AUTOMATED DETECTION OF SATELLITE TRAILS IN}\par
\vspace{0.05em}
{\large\bfseries GROUND-BASED OBSERVATIONS}\par
\vspace{0.25em}
{\normalsize U-Net segmentation with probabilistic Hough post-processing}\par
\vspace{0.3em}
{\small Baron Gracias \,|\, MPhil Data Intensive Science \,|\, University of Cambridge \,|\, Supervisor: Dr Eduardo Gonzalez-Solares}
\end{center}
\vspace{-0.35em}
\noindent\textcolor{rulegray}{\rule{\textwidth}{0.4pt}}
\vspace{0.55em}
]

\section*{Motivation \& Scientific Justification}
Low-Earth-orbit satellite constellations increasingly cross wide-field
astronomical exposures, leaving bright linear trails that corrupt photometry,
astrometry, and transient-source catalogues. Manual inspection does not scale
to survey operation, so observatories need pixel-reliable automated masks that
can run inside reduction pipelines.

This project replicates the satellite-trail detector of Stoppa et al. (2024):
a U-Net trained on labelled MeerLICHT image patches, followed by a probabilistic
Hough transform that reconnects linear trail evidence missed by the network. The
replication is used diagnostically: it tests where the residual error comes
from, not only whether a headline score reproduces.

\section*{Key Findings}
The replication reproduces the target paper's recall closely: five-seed
sampled-test recall is $0.923\pm0.007$ against the paper's reported $0.94$.
Strict precision is lower, at $0.847\pm0.010$ on the sampled test set and
$0.780$ on the parity set that restores the natural negative-patch prevalence.
This shows that precision is prevalence-sensitive, whereas recall is the cleaner
like-for-like comparison across the two test populations. Image-cluster
uncertainty is the relevant interval because source images, not patches, are the
independent units.

\summaryfigure[0.55em]{0.88}{../results/figures/thesis_1_multiseed_replication.pdf}{Five-seed
replication performance for the locked U-Net compared with Stoppa et al.\
Paper-reference markers are shown only for the explicitly reported
precision and recall values; recall is closely reproduced while strict precision
remains lower, motivating the boundary-near error analysis.}

The Hough stage reproduces the paper's intended completeness behaviour. On the
full-coverage tiled canvas, trail-pixel recall rises from $0.918$ to $0.990$,
equivalently reducing the pixel-level false-negative rate from $0.082$ to
$0.0105$. This does not make the Hough mask a better strict pixel classifier:
post-Hough line drawing increases false-positive area, reducing strict precision
from $0.780$ to $0.410$ on the parity canvas.

\summaryfigure[0.55em]{0.88}{../results/figures/hough_gap_bridge_example.pdf}{Hough
bridging a U-Net gap over a bright star:
linear post-processing recovers ground-truth trail pixels missed by the
thresholded U-Net.}

Diagnostics place much of the strict precision gap near annotated boundaries.
False-positive decomposition shows that only $17.5\,\%$ of false-positive
pixels occur in trail-free patches; the remaining $82.5\,\%$ fall within
trail-bearing patches. A one-pixel boundary tolerance raises five-seed precision
from $0.847$ to $0.946$, and $57\,\%$ of all false positives lie within one
pixel of a labelled trail. These diagnostics are consistent with residual error
concentrated near labelled trail boundaries. A blinded single-author
re-annotation of $64$ sealed crops is consistent with this: the original masks,
the re-annotation, and the model all share a median width of $6$\,px, and
the model agrees with the re-annotation about as well as the original labels do,
with no strict-$F_1$ difference resolved under this design (paired $-0.0001$,
$95\%$ CI spanning zero). This is a single-author comparison between mask
sources, not inter-annotator agreement. A small over-firing residual remains on
adversarial controls, so the audit supports a boundary-scale interpretation
without implying a uniformly thin label or model over-paint.

\vspace{0.35em}
\summaryfigure{0.82}{../results/figures/thesis_2_boundary_tolerance.pdf}{A
one-pixel tolerance largely closes the strict precision gap, consistent with
residual disagreement concentrated near labelled boundaries.}

\clearpage

\section*{Methodology}
The study uses 178 labelled MeerLICHT image/mask pairs split at image level into
122 training, 24 validation, and 32 test images, avoiding same-image leakage.
Images are tiled into $528\times528$ patches. The U-Net is reimplemented in
PyTorch from the released Keras model configuration, preserving average pooling,
LeakyReLU activations, in-block dropout, Keras-style initialisation, and a
combined binary cross-entropy plus squared-Dice loss.

Hyperparameters are selected by validation-only Optuna search, the top trials
are retrained across five seeds, and the final threshold is chosen from a
validation precision--recall sweep before test metrics are inspected. Hough
post-processing is evaluated on a full-coverage tiled canvas that reproduces
the parity canvas pixel-for-pixel, with separate precision and completeness
caveats.

\vspace{0.45em}
\summaryfigure{0.88}{../results/figures/predictions/overlay_grid_winner_t44_s2804.png}{Diagnostic
test-patch predictions for the locked U-Net: ground truth in green,
prediction in magenta, and overlap inside the green band.}

\section*{Extensions \& Transfer}
Follow-up experiments test whether the residual can be moved by model or
inference changes. A classifier-gated detector shifts precision by only
$+1.2$ percentage points, consistent with the $17.5\,\%$ inter-patch upper
bound. An attention-gated U-Net at matched capacity is a null result, clDice
shows no hidden topology gain, hard-negative mining adds only $+0.007$
precision, and ensembling with test-time augmentation sits within seed noise.

\vspace{0.45em}
\summaryfigure{0.95}{../results/figures/decam_cold_raw_vs_overlay.pdf}{Cold-domain
DECam examples: raw measured-streak frames and the locked MeerLICHT-trained
U-Net contour. This is a qualitative transfer illustration, not a benchmark.}

\newpage

\section*{Recommendation \& Impact}
The practical result is a high-completeness trail masker whose remaining
strict-precision gap is largely boundary-near. The reporting standard should be
strict precision/recall plus a boundary-tolerant operating
point and an explicit mask-width convention. The strongest remaining check is a
multi-annotator annotation audit, which would establish inter-annotator
agreement and test how much of the strict precision gap reflects mask
granularity rather than detector error. A labelled cross-instrument test is the
next generalisation check; the current DECam application is qualitative only.

The U-Net plus Hough pipeline closely reproduces the target detector's recall
and Hough completeness. The tested capacity, training-protocol and
post-processing changes do not materially reduce the remaining strict-precision
residual. For operational masking, the evidence points to reporting strict
scores alongside boundary-tolerant scores and annotation conventions rather than
prioritising another architecture tweak.

\section*{References}
{\footnotesize
[1] F. Stoppa et al., ``Automated detection of satellite trails in ground-based
observations using U-Net and Hough transform,'' \emph{Astronomy \& Astrophysics},
692, A199, 2024.
[2] O. Ronneberger, P. Fischer, and T. Brox, ``U-Net: Convolutional networks for
biomedical image segmentation,'' MICCAI, 2015.
[3] J. Matas, C. Galambos, and J. Kittler, ``Robust detection of lines using the
progressive probabilistic Hough transform,'' \emph{Computer Vision and Image
Understanding}, 78(1), 2000.
}

\section*{Declaration of Use of Autogeneration Tools}
{\footnotesize
I have used Codex (ChatGPT) and Claude Code to support me across this work,
primarily in the following areas:
\begin{itemize}
  \item Providing advice and usage support for libraries like PyTorch and PIL.
  \item Generating boilerplate code for plotting, testing and scripting.
  \item Generating function docstrings.
  \item Reviewing and auditing project structure for completeness, consistency,
        accuracy and best practices in software engineering.
\end{itemize}
}

\end{document}
